\subsection{Injective paths}

PROPOSITION 18. --- Let X be a Hausdorff topological space. Let a and b be distinct points of X that belong to the same path-connected component of X. There exists an injective path connecting a to b in X.

Let \(f \colon \mathbf{I} \to \mathbf{X}\) be a path connecting \(a\) to \(b\) in \(\mathbf{X}\) . Let \(\mathcal{U}\) be the set of open subsets \(\mathrm{U}\) of {]}0,1{[} such that \(f(x) = f(y)\) for every connected component {]}x, y{[} of U.

Lemma 8. --- The set U, ordered by inclusion, is inductive.

Let V be a totally ordered subset of U. Let us prove that the union V of the sets belonging to V is an element of U, that is, that, for every connected component {]}u, v{[} of V, we have \(f(u) = f(v)\) .

Let \(x\) be a point of \([u, v[\) . Let \(\mathcal{V}_x\) be the set of \(U \in \mathcal{V}\) such that \(x \in U\) ; for such a \(U\) , denote by \([u_U, v_U[\) the connected component of \(x\) in \(U\) . We have \(u_{U'} \leqslant u_U < v_U \leqslant v_{U'}\) if \(U\) and \(U'\) are elements of \(\mathcal{V}_x\) such that \(U \subset U'\) . Since the union over \(U \in \mathcal{V}_x\) of the \([u_U, v_U[\) is equal to {]}u, v{[} by Lemma 7 of III, p.~280, we have u = lim \(u_{U}\) and v = lim \(v_{U}\) , where the limits are taken over the directed set \(V_{x}\) . Since the map f is continuous and the topological space X is Hausdorff, the equalities \(f(u_{\mathrm{U}}) = f(v_{\mathrm{U}})\) for all \(U \in V_{x}\) imply that \(f(u) = f(v)\) , which is what had to be proved.

By virtue of E, III, p.~20, th. 2, there exists an open subset U belonging to \(\mathcal{U}\) which is maximal for the inclusion relation.

Let \(g \colon \mathbf{I} \to \mathrm{X}\) be the map defined in the following way. If \(t \notin \mathrm{U}\) we set \(g(t) = f(t)\) ; otherwise, denoting by \(]u, v[\) the connected component of \(t\) in \(\mathrm{U}\) we set \(g(t) = f(u) = f(v)\) , so that the map \(g|[u, v]\) is constant with image \(f(u)\) .

Prop. 18 follows from the following lemma.

Lemma 9. --- There exists a continuous, increasing, and surjective map \(u\colon\mathbf{I}\to\mathbf{I}\) and an injective path \(c\colon\mathbf{I}\to\mathbf{X}\) connecting \(a\) to \(b\) such that \(g=c\circ u\) .

Let us first prove that the map \(g\) is continuous. Let \(t\) be a point of \(\mathbf{I}\) . If \(t \in \mathrm{U}\) , \(g\) is constant in a neighborhood of \(t\) , hence continuous at \(t\) . Suppose \(t \notin \mathrm{U}\) and let \(\mathrm{W}\) be an open neighborhood of \(g(t)\) in \(\mathrm{X}\) . Let \(\mathrm{V}\) be an open interval in \(\mathbf{I}\) containing \(t\) such that \(f(\mathrm{V}) \subset \mathrm{W}\) ; such an interval exists since \(f\) is continuous at \(t\) . Let us prove that \(g(\mathrm{V}) \subset \mathrm{W}\) . Let \(x \in \mathrm{V}\) ; if \(x \notin \mathrm{U}\) , we have \(g(x) = f(x)\) , therefore \(g(x) \in \mathrm{W}\) . Suppose then \(x \in \mathrm{U}\) and let \(]u, v[\) be the connected component of \(x\) in \(\mathrm{U}\) . Observe that \(]u, v[\) does not contain \(t\) . Consequently, if \(x < t\) , we have \(x < v \leqslant t\) , whence \(v \in \mathrm{V}\) and \(g(x) = g(v) = f(v) \in f(\mathrm{V}) \subset \mathrm{W}\) . One proves similarly that \(g(x) \in \mathrm{W}\) if \(x > t\) . Thus, \(g\) is continuous at \(t\) .

Let \(u\) and \(v\) be points of \(\mathbf{I}\) such that \(g(u) = g(v)\) and \(u < v\) . Let us prove that \(]u, v[ \subset U\) . Set \(U' = U \cup ]u, v[\) ; it is an open subset of \(]0, 1[\) .

If \(u\) belongs to U, denote by \(u'\) the greatest lower bound in I of the connected component of \(u\) in U; set \(u' = u\) if \(u\) does not belong to U. Likewise, if \(v\) belongs to U, denote by \(v'\) the least upper bound in I of the connected component of \(v\) in U; set \(v' = v\) if \(v\) does not belong to U. Then, \(]u', v'[\) is a connected component of U', and we have \(f(u') = g(u) = g(v) = f(v')\) . Since the connected components of U' distinct from \(]u', v'[\) are connected components of U, the open set U' is an element of \(\mathcal{U}\) . Since U is a maximal element of \(\mathcal{U}\) for the inclusion relation, we have U' = U and \(]u, v[\) is contained in U. This proves that \(g\) is constant on the interval \([u, v]\) . The fibers of \(g\) are therefore intervals of I, and these intervals are closed because g is continuous.

Let R denote the equivalence relation associated with g, and let p be the canonical surjection from I onto I/R. There exists a unique continuous map \(g'\) from I/R to X such that \(g = g' \circ p\) ; this map is injective. By the corollary, III, p.~276, of Prop. 16, there exists a map u, increasing, continuous, and surjective, such that R is the equivalence relation associated with u. Since the space I is compact and the space X is Hausdorff, the map u is closed, hence strict (I, p.~18, Example 2). It defines, by passing to the quotient, a homeomorphism \(u'\) from I/R onto I. Then, the map \(g' \circ (u')^{-1}\) from I into X is an injective path with initial point a and terminal point b.

